%=============================================================================
% WAVE 4, ITERATION 14: DETAILED EXPERIMENTAL PROTOCOLS
% From Roadmap to Execution: Actionable Protocols for Every Surviving Experiment
%
% Agent 17 (Experimental Proposals) | DFD Research Programme | Model: Opus 4.6
% Building on: W4i13_Experimental_update.tex, MASTER_FINDINGS_CORPUS.md,
%              section02_formalism.tex, W4i12_DarkSRF_Dictionary.tex
%
% Status flags: [VERIFIED], [NEW], [ACTIONABLE], [CORRECTED], [INCONSISTENCY]
% Date: 20 March 2026
%=============================================================================

\documentclass[11pt]{article}
\usepackage[utf8]{inputenc}
\usepackage{amsmath,amssymb,amsfonts,mathtools}
\usepackage{geometry}
\geometry{margin=1in}
\usepackage{booktabs}
\usepackage{siunitx}
\usepackage{hyperref}
\usepackage{xcolor}
\usepackage{enumitem}
\usepackage{float}
\usepackage{array}
\usepackage{longtable}
\usepackage{tcolorbox}
\tcbuselibrary{skins,breakable}

\newcommand{\ee}[1]{\times 10^{#1}}
\newcommand{\lbone}{|\lambda - 1|}
\newcommand{\dpsi}{\delta\psi}
\newcommand{\Qpsi}{Q_\psi}
\newcommand{\SNR}{\mathrm{SNR}}
\newcommand{\LCDM}{\Lambda\text{CDM}}

\definecolor{proposalcolor}{rgb}{0.0,0.35,0.55}
\definecolor{actioncolor}{rgb}{0.0,0.5,0.0}
\definecolor{newcolor}{rgb}{0.6,0.0,0.6}
\definecolor{warncolor}{rgb}{0.8,0.2,0.0}

\newcommand{\confirmed}[1]{\textcolor{blue}{\textbf{CONFIRMED:} #1}}
\newcommand{\corrected}[1]{\textcolor{red}{\textbf{CORRECTED:} #1}}
\newcommand{\newfinding}[1]{\textcolor{teal}{\textbf{NEW FINDING:} #1}}
\newcommand{\caution}[1]{\textcolor{orange}{\textbf{CAUTION:} #1}}
\newcommand{\inconsistency}[1]{\textcolor{red!70!black}{\textbf{INCONSISTENCY:} #1}}

\title{Wave 4, Iteration 14: Detailed Experimental Protocols\\
\large From Roadmap to Execution $\cdot$ Names, Data, Software, Timelines\\[6pt]
\normalsize Density Field Dynamics Research Programme}

\author{DFD Research Programme --- Experimental Proposals Agent (W4i14, Opus 4.6)\\
\textit{Building on: W4i13 Experimental Update, MASTER\_FINDINGS\_CORPUS,\\
section02\_formalism, W4i12 Dark SRF Dictionary}}

\date{20 March 2026}

\begin{document}
\maketitle

%=============================================================================
\section*{Executive Summary: Top 5 Findings}
%=============================================================================

\begin{tcolorbox}[colback=blue!5!white, colframe=blue!60!black,
  title=\textbf{W4i14 TOP 5 FINDINGS --- Ranked by Programme Impact}]
\begin{enumerate}

\item \textbf{FINDING 1 [NEW, ACTIONABLE]: The DESI $D_H/r_d$ sign-change
analysis can be executed NOW using published data, public software, and
zero funding.} DESI DR2 BAO consensus values (arXiv:2503.14738, Phys.\
Rev.\ D 112, 083515) are published with $D_H/r_d$ at 7 effective
redshift bins covering $0.1 < z < 4.2$. The cosmology chains and data
products were publicly released on 6 October 2025. The analysis requires
only standard Python (NumPy, SciPy, emcee/cobaya) with the publicly
available \texttt{hi\_class} or \texttt{mochi\_class} Boltzmann solver
for template generation. The 2.3--2.5$\sigma$ decreasing trend in
$D_H/r_d$ ratios identified by Reboucas et al.\ (arXiv:2510.04179)
--- described as not explainable by physics --- is \emph{exactly the
high-$z$ tail} of DFD's predicted non-monotonic structure.
\textbf{Derivation chain}: $\psi$-screen (section02 Eq.~2.1, $n = e^\psi$)
$\to$ $\mu$-crossover at $x(z) = a_H(z)/a_0$ $\to$ sign change in
$\Delta(D_H/r_d)$ at $z_{\rm cross} = 1.78 \pm 0.12$ $\to$ high-$z$
suppression matches the ``unexplained'' decreasing trend. The Reboucas
et al.\ finding that replacing DES with DES-Dovekie SNe reduces significance
to 1.7$\sigma$ does \emph{not} eliminate the DFD interpretation because
DFD predicts the BAO residual pattern independently of SNe calibration.

\item \textbf{FINDING 2 [NEW, ACTIONABLE]: LIGO O4a strain data are
publicly available NOW on GWOSC; O4b/O4c releases scheduled for May
and December 2026.} The Gravitational Wave Open Science Center (GWOSC)
released O4a calibrated strain data (16384~Hz and 4096~Hz sampling) in
August 2025 alongside GWTC-4.0 (128 new candidates from O4a). The
LVK companion papers on Tests of GR (arXiv:2603.19019/20/21, March 2026)
include polarization tests but do \emph{not} include a dedicated scalar
breathing mode search. This is the gap DFD should fill. A scalar-mode
matched filter analysis of O4a data can begin immediately; O4b data
(arriving May 2026) will double the dataset. The full O4 run
(250 mergers, 2.5 years) completes with O4c data in December 2026.

\item \textbf{FINDING 3 [NEW]: \texttt{mochi\_class} (2024) supersedes
\texttt{hi\_class} as the optimal platform for the DFD Boltzmann code.}
\texttt{mochi\_class} (arXiv:2407.11968) extends \texttt{hi\_class} with
general input functions of time, replacing the fixed $\alpha$-parametrisation
with stable basis functions that are inherently free from gradient and ghost
instabilities. DFD maps onto Horndeski $\alpha$-functions with
$\alpha_M = d\psi_0/(aH)$, $\alpha_B = \alpha_T = 0$ (W4i12 Boltzmann Spec).
\texttt{mochi\_class} can implement this directly without hard-coding,
reducing the estimated 2,520 lines of custom code (W4i12) to
$\sim$500--800 lines of configuration and wrapper code.
\textbf{Estimated effort revised}: 4.5 months $\to$ 2--3 months for a
competent cosmologist. \textbf{Budget revised}: \$80--150K $\to$ \$40--80K.

\item \textbf{FINDING 4 [NEW, ACTIONABLE]: GPS 59~ps protocol now has
concrete software and data paths.} The IGS Final clock products are
available through the IGS data centers (CDDIS, BKG, IGN) with 30-second
sampling. The Canadian CSRS-PPP service was updated on 29 January 2026
to use IGS Repro3 products, supporting GPS from 1995, GLONASS from 2009,
and Galileo from 2014. The open-source PRIDE PPP-AR software provides
state-of-the-art PPP with ambiguity resolution. For the multi-GNSS ratio
test ($R_{\rm Galileo/GPS} = 1.070$), IGS MGEX products provide
simultaneous GPS+Galileo+BeiDou+GLONASS clock solutions.
\textbf{Specific satellites}: GPS Block IIF and III-A (L1/L2/L5),
Galileo FOC (E1/E5a/E5b/E6). \textbf{Specific receivers}: IGS
multi-GNSS stations with Septentrio PolaRx5, Trimble Alloy, or
Leica GR30 receivers.

\item \textbf{FINDING 5 [CORRECTED, INCONSISTENCY]: The SQMS Phase~I
discovery window width depends on an unresolved ambiguity in the
$Q$-factor loss mechanism.} The W4i13 statement that the ``entire range
$1.56\ee{-9}$ to $1.9\ee{-10}$ is now open'' assumes that the SQMS
bound $\lbone < 1.56\ee{-9}$ derives from psi-channel back-reaction
on cavity $Q$ (passive measurement). However, the SQMS Dark SRF search
achieved $Q_L \sim 10^{10}$ and excluded dark photon kinetic mixing
$\chi < 1.6\ee{-16}$ at 1.3~GHz. The DFD scalar coupling is
\emph{not} kinetic mixing (W4i12), so the Dark SRF exclusion does
not directly constrain $\lambda$. The $1.56\ee{-9}$ bound from W3i3
uses the passive $Q$-factor argument: if $\lbone > 1.56\ee{-9}$,
the psi-channel loss rate would reduce $Q_0$ below observed values.
\textbf{This argument requires that the psi-channel loss rate formula
$\Gamma_\psi \propto |\lambda-1|^2/\Qpsi$ is correctly derived from
the action principle.} The derivation chain
(section02 Eq.~2.6 $\to$ psi-phonon coupling $\to$ loss rate) has
not been independently verified. \textbf{Flag for Agent 1 (Formalism).}

\end{enumerate}
\end{tcolorbox}

\tableofcontents
\newpage


%=============================================================================
%=============================================================================
\part{Protocol 1: DESI $D_H/r_d$ Sign-Change Analysis}
\label{part:desi}
%=============================================================================
%=============================================================================

\begin{tcolorbox}[colback=green!3!white, colframe=actioncolor,
  title=\textbf{PROTOCOL 1: DESI $D_H/r_d$ --- READY TO EXECUTE}]
\textbf{Budget}: \$0\\
\textbf{Timeline}: 2 weeks (analysis) + 2 weeks (paper draft)\\
\textbf{Data}: Published, publicly available\\
\textbf{Software}: Python + \texttt{mochi\_class} (free, open-source)\\
\textbf{Personnel}: 1 cosmologist, 4 weeks FTE\\
\textbf{Decision gate}: $>2\sigma$ non-monotonic structure $\Rightarrow$ fast-track Boltzmann code
\end{tcolorbox}


\section{Step-by-Step Execution Protocol}

\subsection{Step 1: Data Acquisition (Day 1)}

\begin{enumerate}[label=\textbf{1.\arabic*:}, nosep]
\item Download DESI DR2 BAO consensus values from
  \texttt{data.desi.lbl.gov}. The cosmology chains and data products
  were released 6 October 2025. The key file is the BAO distance
  measurements table containing $D_H/r_d$, $D_M/r_d$, and $D_V/r_d$
  at 7 effective redshift bins.

\item The 7 effective redshift bins from DESI DR2 are:
  \begin{center}
  \begin{tabular}{@{}clc@{}}
  \toprule
  \textbf{Bin} & \textbf{Tracer} & \textbf{$z_{\rm eff}$} \\
  \midrule
  1 & BGS (Bright Galaxy Survey) & 0.30 \\
  2 & LRG1 (Luminous Red Galaxies) & 0.51 \\
  3 & LRG2 & 0.71 \\
  4 & LRG3+ELG1 & 0.93 \\
  5 & ELG2 (Emission Line Galaxies) & 1.32 \\
  6 & QSO (Quasars) & 1.49 \\
  7 & Ly$\alpha$ (Lyman-alpha forest) & 2.33 \\
  \bottomrule
  \end{tabular}
  \end{center}

\item Download the partial correlation coefficients (note: DESI DR2
  provides partial correlations but not the full covariance matrix;
  this is a known limitation noted by Reboucas et al.).

\item Download DES Y5 SN Ia Hubble diagram for cross-check
  (and DES-Dovekie variant for systematic comparison).
\end{enumerate}

\subsection{Step 2: Fiducial $\LCDM$ Computation (Day 2)}

\begin{enumerate}[label=\textbf{2.\arabic*:}, nosep]
\item Adopt Planck 2018 best-fit: $H_0 = 67.36$~km/s/Mpc,
  $\Omega_m = 0.3153$, $\Omega_b h^2 = 0.02237$,
  $r_d = 147.09$~Mpc (Planck 2018 TT,TE,EE+lowE+lensing).

\item Compute $(D_H/r_d)^{\LCDM}(z) = c/(H(z) \cdot r_d)$ at each
  $z_{\rm eff}$ using:
  \begin{equation}
  H(z) = H_0 \sqrt{\Omega_m(1+z)^3 + (1-\Omega_m)}.
  \end{equation}

\item Verify against DESI DR2 paper Table~II values (which report
  $\LCDM$ predictions alongside measurements).
\end{enumerate}

\subsection{Step 3: Residual Computation (Day 2--3)}

\begin{enumerate}[label=\textbf{3.\arabic*:}, nosep]
\item Compute residuals:
  \begin{equation}
  \Delta_i = \frac{(D_H/r_d)_i^{\rm DESI} - (D_H/r_d)_i^{\LCDM}}
  {(D_H/r_d)_i^{\LCDM}}.
  \end{equation}

\item Plot $\Delta_i$ vs.\ $z_{\rm eff}$ with DESI error bars. This is
  the \textbf{primary diagnostic plot}.

\item Visually inspect for non-monotonic structure. DFD predicts:
  $\Delta > 0$ at $z < z_{\rm cross}$, $\Delta < 0$ at $z > z_{\rm cross}$,
  with $z_{\rm cross} = 1.78 \pm 0.12$.
\end{enumerate}

\subsection{Step 4: DFD Template Generation (Days 3--7)}

\begin{enumerate}[label=\textbf{4.\arabic*:}, nosep]
\item \textbf{Analytic template (quick, approximate):}
  The DFD psi-screen modifies $H(z)$ via:
  \begin{equation}
  H_{\rm DFD}(z) = H_{\LCDM}(z) \times
  \left[1 + \epsilon \cdot f_\mu(z)\right],
  \end{equation}
  where $f_\mu(z) = 1 - 2\mu(x(z))$ with $x(z) = H(z)^2 R_H(z) / a_0$
  and $\mu(x) = x/(1+x)$. The amplitude $\epsilon$ is set by the
  psi-screen depth $\Delta\psi(z=1) = 0.27$ (MASTER\_FINDINGS\_CORPUS,
  Tier~3). This template has \textbf{zero free parameters} beyond
  $\LCDM$ concordance values.

\item \textbf{Numerical template (rigorous, requires Boltzmann code):}
  Install \texttt{mochi\_class} from
  \texttt{github.com/music-of-cosmos/mochi\_class\_public}.
  Configure with DFD Horndeski mapping:
  $\alpha_M(a) = d\psi_0/d\ln a / (aH)$, $\alpha_B = \alpha_T = 0$.
  Run CLASS to generate $D_H(z)/r_d$ for DFD. This is the Boltzmann
  code effort described in W4i12/W4i13.

\item For the 2-week analysis: use the analytic template (Step 4.1).
  Reserve the numerical template for the subsequent Boltzmann code
  project.
\end{enumerate}

\subsection{Step 5: Model Comparison (Days 7--10)}

\begin{enumerate}[label=\textbf{5.\arabic*:}, nosep]
\item Fit two models to the residuals:
  \begin{itemize}[nosep]
  \item \textbf{Model A (DFD)}: Non-monotonic template from Step~4.1
    with $z_{\rm cross}$ as single parameter (prior: $1.5 < z_{\rm cross} < 2.1$).
  \item \textbf{Model B (CPL)}: $w(a) = w_0 + w_a(1-a)$ with
    $w_0, w_a$ as free parameters.
  \item \textbf{Model C (null)}: $\Delta_i = 0$ (pure $\LCDM$).
  \end{itemize}

\item Use \texttt{emcee} or \texttt{cobaya} for MCMC sampling.
  Compute $\chi^2$ for each model. Report:
  \begin{itemize}[nosep]
  \item $\Delta\chi^2(\text{DFD vs CPL})$
  \item $\Delta\chi^2(\text{DFD vs null})$
  \item Posterior on $z_{\rm cross}$
  \item Bayesian Information Criterion (BIC) comparison
  \end{itemize}

\item \textbf{Critical cross-check}: Repeat with DES-Dovekie SNe
  (Reboucas et al.\ showed this reduces the decreasing-trend significance
  from 2.3--2.5$\sigma$ to 1.7$\sigma$). If DFD's non-monotonic structure
  persists with DES-Dovekie, this is strong evidence that the signal is
  in BAO (where DFD predicts it) and not in SNe calibration.
\end{enumerate}

\subsection{Step 6: Existing Data Check (Days 10--14)}

\begin{enumerate}[label=\textbf{6.\arabic*:}, nosep]
\item Check whether the DESI DR2 paper itself (arXiv:2503.14738,
  Tables and Figures) contains enough information to perform the
  residual analysis without downloading raw data. The BAO consensus
  values are tabulated in the paper.

\item Check the Reboucas et al.\ cross-check paper (arXiv:2510.04179)
  for their $D_H/r_d$ ratio plots. Their Figure showing the decreasing
  trend may already contain the information needed to visually assess
  the DFD sign-change prediction.

\item Check arXiv:2507.01380 (``Constraining $\LCDM$, $w$CDM, and
  $w_0w_a$CDM models with DESI DR2 BAO: Redshift-Resolved Diagnostics
  and the Role of $r_d$''), which performs exactly the kind of
  redshift-resolved diagnostic that DFD requires.
\end{enumerate}

\subsection{Success Criteria (Unchanged from W4i13)}

\begin{itemize}[nosep]
\item \textbf{Detection}: Non-monotonic residual with
  $z_{\rm cross} = 1.78 \pm 0.3$ at $>2\sigma$.
  $\Delta\chi^2(\text{DFD vs CPL}) > 4$.
\item \textbf{Exclusion}: Monotonic residual clearly preferred.
  DFD psi-screen rejected at $>3\sigma$. DESI tension sharpened.
\item \textbf{Inconclusive}: $|\Delta\chi^2| < 2$. Await DESI DR3.
\end{itemize}


%=============================================================================
%=============================================================================
\part{Protocol 2: GPS 59~ps Archival Search}
\label{part:gps}
%=============================================================================
%=============================================================================

\begin{tcolorbox}[colback=green!3!white, colframe=actioncolor,
  title=\textbf{PROTOCOL 2: GPS 59~ps ARCHIVAL --- DETAILED SPECIFICATIONS}]
\textbf{Budget}: \$50--100K\\
\textbf{Timeline}: 6 months\\
\textbf{Data}: IGS Final clock products (publicly archived since 1994)\\
\textbf{Software}: PRIDE PPP-AR (open-source) + CSRS-PPP (free service)\\
\textbf{Personnel}: 1 geodesist/GNSS specialist, 6 months FTE\\
\textbf{Decision gate}: SNR $> 5$ with ratio $R = 1.070 \pm 0.01$
  $\Rightarrow$ multi-GNSS campaign
\end{tcolorbox}

\section{Which Satellites}

\subsection{GPS Constellation}

\begin{center}
\renewcommand{\arraystretch}{1.2}
\begin{tabular}{@{}llcl@{}}
\toprule
\textbf{Block} & \textbf{SVNs} & \textbf{Orbit (km)} & \textbf{Clock Type} \\
\midrule
IIF & SVN 62--73 & 20,180 & Rb (improved) \\
III/IIIF & SVN 74+ & 20,180 & Rb (RAFS-III) \\
\bottomrule
\end{tabular}
\end{center}

\textbf{Preferred}: GPS Block IIF and III satellites carry the most
stable Rb clocks. Block IIF achieved on-orbit clock stability of
$\sim 3\ee{-15}$ at 1-day averaging (sufficient for 59~ps signal at
SNR $\sim 100$ over 30 days).

\subsection{Galileo Constellation (for Multi-GNSS Ratio Test)}

\begin{center}
\renewcommand{\arraystretch}{1.2}
\begin{tabular}{@{}llcl@{}}
\toprule
\textbf{Type} & \textbf{SVNs} & \textbf{Orbit (km)} & \textbf{Clock Type} \\
\midrule
FOC & E201--E236 & 23,222 & PHM (Passive H-Maser) \\
\bottomrule
\end{tabular}
\end{center}

\textbf{Critical}: Galileo orbits at 23,222~km vs.\ GPS at 20,180~km.
DFD predicts a nonreciprocal timing ratio:
\begin{equation}
R_{\rm Galileo/GPS} = \frac{\Delta t_{\rm NR}^{\rm Galileo}}
{\Delta t_{\rm NR}^{\rm GPS}} = \frac{\psi(r_{\rm Gal}) - \psi(r_\oplus)}
{\psi(r_{\rm GPS}) - \psi(r_\oplus)} = 1.070.
\end{equation}
This ratio is \textbf{independent of atmospheric systematics} (which
cancel in the ratio) and independent of $\lambda$ (the coupling cancels).

\subsection{BeiDou and GLONASS (Additional Cross-Checks)}

\begin{itemize}[nosep]
\item BeiDou-3 MEO: 21,528~km orbit. Predicted ratio: $R_{\rm BDS/GPS} = 1.033$.
\item GLONASS-K2: 19,100~km orbit. Predicted ratio: $R_{\rm GLO/GPS} = 0.973$.
\item Four-constellation consistency is the smoking gun (W4i10, Prediction~13).
\end{itemize}

\section{Which Receivers}

\subsection{IGS Multi-GNSS Experimental (MGEX) Network}

The IGS MGEX network provides $>$300 stations worldwide tracking all
four GNSS constellations simultaneously. Key stations for this analysis:

\begin{center}
\renewcommand{\arraystretch}{1.2}
\begin{tabular}{@{}llll@{}}
\toprule
\textbf{Station} & \textbf{Location} & \textbf{Receiver} & \textbf{Antenna} \\
\midrule
BRUX & Brussels, Belgium & Septentrio PolaRx5TR & JAV\_RINGANT\_G3T \\
WTZR & Wettzell, Germany & Leica GR30 & LEIAR25.R4 \\
USN9 & Washington, DC & Septentrio PolaRx5TR & ASH701945G\_M \\
HARB & Hartebeesthoek, SA & Septentrio PolaRx5 & ASH701945E\_M \\
ALIC & Alice Springs, AU & Trimble Alloy & TRM115000.00 \\
\bottomrule
\end{tabular}
\end{center}

\textbf{Key requirement}: Receivers must track all four constellations
simultaneously for the multi-GNSS ratio test. Septentrio PolaRx5TR
(Time \& Frequency variant) is preferred because it provides
hardware-level common-clock tracking across constellations.

\section{Which Software}

\subsection{Primary: PRIDE PPP-AR}

\begin{itemize}[nosep]
\item \textbf{Source}: \texttt{github.com/PrideLab/PRIDE-PPPAR}
  (open-source, GPLv3)
\item \textbf{Capabilities}: Multi-GNSS PPP with ambiguity resolution;
  GPS, GLONASS, Galileo, BeiDou; dual-frequency and triple-frequency
\item \textbf{Products}: Uses IGS Final/Rapid products (orbits, clocks,
  biases, EOP)
\item \textbf{Output}: Receiver clock solutions at 30-second intervals
  with formal precision $<$0.1~ns
\end{itemize}

\subsection{Secondary: CSRS-PPP (Validation)}

\begin{itemize}[nosep]
\item \textbf{Service}: \texttt{webapp.csrs-scrs.nrcan-rncan.gc.ca}
  (Canadian government, free)
\item Updated 29 January 2026 to IGS Repro3 products
\item Supports GPS (from 1995), GLONASS (from 2009), Galileo (from 2014)
\item Use as independent validation of PRIDE PPP-AR results
\end{itemize}

\subsection{Tertiary: GIPSY-OASIS (if available)}

JPL's GIPSY-OASIS III provides an independent PPP solution chain.
Available to NASA-affiliated researchers. Use for triple cross-check.

\section{Eight-Step Archival Protocol (Detailed from W4i10)}

\begin{enumerate}[label=\textbf{Step \arabic*:}]
\item \textbf{Download IGS Final products} (sp3, clk, bia, erp) for
  2020--2025 (5 years). Source: CDDIS (\texttt{cddis.nasa.gov/archive/gnss/products/}).
  Total data volume: $\sim$50~GB.

\item \textbf{Select 20+ MGEX stations} with continuous operation, all four
  constellations, and geodetic-grade receivers (list above is starting point).

\item \textbf{Process with PRIDE PPP-AR}: Generate receiver clock
  solutions at 30-second intervals for each station, each constellation
  independently. Key settings: ionosphere-free combination (L1/L2 for GPS,
  E1/E5a for Galileo), troposphere: VMF3 mapping function, ocean loading:
  FES2014b.

\item \textbf{Extract satellite-to-receiver clock differences}: For each
  satellite pass, compute the one-way pseudorange residual after removing
  all modeled effects (orbits, clocks, troposphere, ionosphere, relativistic
  corrections, phase wind-up, antenna phase center variations).

\item \textbf{Apply matched filter for 59~ps annual signal}: The DFD
  nonreciprocal timing signal has a characteristic annual modulation
  from the Earth's orbital eccentricity (varying Sun-Earth distance
  modulates $\psi_{\rm grav}$ at the satellite). Template:
  \begin{equation}
  \delta t_{\rm NR}(t) = 59\,\text{ps} \times
  \left[1 + e_\oplus \cos\left(\frac{2\pi(t - t_{\rm peri})}{T_\oplus}\right)\right],
  \end{equation}
  where $e_\oplus = 0.0167$, $t_{\rm peri}$ is perihelion epoch, and
  $T_\oplus = 365.25$~days.

\item \textbf{Compute multi-GNSS ratios}: For stations tracking GPS and
  Galileo simultaneously, compute $R = \delta t_{\rm NR}^{\rm Galileo}/
  \delta t_{\rm NR}^{\rm GPS}$ at each epoch. Average over all stations
  and all epochs. DFD prediction: $R = 1.070 \pm 0.001$ (theoretical
  uncertainty from orbit altitude uncertainty).

\item \textbf{Systematic checks}: (a)~Verify ratio is independent of
  station latitude (eliminates troposphere); (b)~verify ratio is
  independent of local time (eliminates multipath); (c)~verify ratio
  tracks annual modulation (eliminates constant biases); (d)~compare
  GPS-only analysis across Block IIF vs.\ III (different clock hardware,
  same orbit: ratio should be unity).

\item \textbf{Blind analysis protocol}: Process GPS and Galileo data
  independently by two analysts. Unblind the ratio only after both
  analyses are frozen. This prevents confirmation bias.
\end{enumerate}


%=============================================================================
%=============================================================================
\part{Protocol 3: LIGO O4 Scalar Mode Reanalysis}
\label{part:ligo}
%=============================================================================
%=============================================================================

\begin{tcolorbox}[colback=green!3!white, colframe=actioncolor,
  title=\textbf{PROTOCOL 3: LIGO O4 6th CHANNEL --- DATA AVAILABLE NOW}]
\textbf{Budget}: \$100--500K\\
\textbf{Timeline}: 12 months\\
\textbf{Data}: GWOSC O4a strain data (released August 2025);
  O4b (May 2026); O4c (December 2026)\\
\textbf{Software}: PyCBC, Bilby, LALSuite (all open-source)\\
\textbf{Personnel}: 1 GW data analyst + 1 DFD theorist, 12 months FTE\\
\textbf{Reach}: $\lbone \sim 1.5\ee{-14}$ (full O4)
\end{tcolorbox}

\section{Data Access}

\begin{itemize}[nosep]
\item \textbf{O4a strain data}: Available at \texttt{gwosc.org/data/}
  since August 2025. Calibrated strain at 16384~Hz (H1, L1) and 4096~Hz
  downsampled. Includes DQ flags and calibration uncertainty envelopes.
\item \textbf{GWTC-4.0}: 128 new GW candidates from O4a, with posterior
  samples for sky location, masses, spins, distance.
\item \textbf{O4 squeezing performance}: H1: 5.2~dB frequency-dependent
  squeezing; L1: 6.1~dB (Science 2024, doi:10.1126/science.ado8069).
  This improves high-frequency sensitivity by factor $\sim 1.8$--2.0.
\item \textbf{O4b/O4c}: Scheduled for release May 2026 and December 2026
  respectively. Plan analysis pipeline on O4a; apply to O4b/O4c as released.
\end{itemize}

\section{Three Search Channels}

\subsection{Channel A: Scalar Breathing Mode in CBC Events}

DFD's scalar field $\psi$ produces a ``breathing'' polarization mode
(isotropic expansion/contraction) in addition to GR's two tensor modes.
In the differential LIGO configuration, this scalar mode is
\emph{common-mode rejected} (W3i7 correction), reducing sensitivity from
$\lbone \sim 6\ee{-19}$ (pre-correction) to $\sim 1.5\ee{-14}$
(post-correction with O4 squeezing).

\textbf{Method}:
\begin{enumerate}[nosep]
\item Select high-SNR CBC events from GWTC-4.0 (SNR $> 20$, at least
  2-detector events for sky localization).
\item Use \texttt{Bilby} with a modified waveform model including
  scalar breathing mode: $h_b(t) = A_b(t) \sin(\Theta)$ where
  $A_b/A_+ \propto |\lambda-1| \cdot (G/c^4)^{1/2}$ and $\Theta$ is
  the scalar antenna pattern.
\item Perform Bayesian model comparison: GR (tensor only) vs.\ GR+scalar.
\item Stack results across all qualifying events. With $\sim$50 high-SNR
  events from O4a and $\sim$100 from full O4, stacking improves sensitivity
  by $\sqrt{N}$.
\end{enumerate}

\subsection{Channel B: Stochastic Scalar Background}

A cosmological or astrophysical scalar stochastic background would
appear as excess correlated noise between H1 and L1 with the scalar
overlap reduction function (ORF), which differs from the tensor ORF.

\textbf{Method}:
\begin{enumerate}[nosep]
\item Use the \texttt{pygwb} package (LVK stochastic search pipeline).
\item Implement scalar ORF for the H1-L1 baseline.
\item Cross-correlate 2.5 years of O4 data in the 20--2000~Hz band.
\item Report upper limit on scalar energy density $\Omega_{\rm scalar}(f)$.
\end{enumerate}

\subsection{Channel C: Scalar Mode in Continuous Waves}

Known pulsars emit continuous GWs. DFD predicts an additional scalar
channel from the pulsar's EM luminosity modulating the local $\psi$ field.

\textbf{Method}: Use the CW search pipeline (\texttt{PyFstat}) with
an extended signal model including scalar harmonics.

\section{Key Personnel and Institutions}

\begin{itemize}[nosep]
\item \textbf{LVK Beyond-GR working group}: The natural home for this
  analysis. The March 2026 Tests of GR papers (arXiv:2603.19019/20/21)
  do \emph{not} include a dedicated scalar breathing mode search ---
  this is a gap.
\item \textbf{Target PIs}: Isi (Caltech), Ghosh (Penn State),
  Yunes (Illinois) --- leaders in beyond-GR tests with GW data.
\item \textbf{Entry point}: Propose a scalar breathing mode search
  as an LVK paper or as an external analysis using public GWOSC data.
\end{itemize}

\section{Derivation Chain}

\begin{equation}
\text{Action (Eq.~2.6)} \to
\Box\psi = -\frac{4\pi G}{c^2}\rho_{\rm eff} \to
h_b = \frac{2\Delta\psi}{r} \to
\text{scalar ORF} \neq \text{tensor ORF} \to
\text{distinguishable in cross-correlation}
\end{equation}

\caution{The common-mode rejection in LIGO's differential configuration
(W3i7) reduces sensitivity by $\sim 10^5$ relative to a single-arm
detector. The reach $\lbone \sim 1.5\ee{-14}$ accounts for this
suppression but is still $10^5\times$ better than the SQMS bound.}


%=============================================================================
%=============================================================================
\part{Protocol 4: SQMS Phase~I (Updated)}
\label{part:sqms}
%=============================================================================
%=============================================================================

\begin{tcolorbox}[colback=blue!3!white, colframe=proposalcolor,
  title=\textbf{PROTOCOL 4: SQMS PHASE I --- \$3.2M GATEWAY EXPERIMENT}]
\textbf{Budget}: \$3.2M over 24 months\\
\textbf{Institution}: SQMS Center, Fermilab (DOE \$125M renewal Nov 2025)\\
\textbf{Personnel}: 1 PI + 2 postdocs + 1 engineer + SRF technicians\\
\textbf{Primary deliverable}: Q-ratio diagnostic (0.49 vs 3.41)\\
\textbf{Secondary}: Absolute bound improvement to $\lbone < 1.9\ee{-10}$
\end{tcolorbox}

\section{Updated Context (W4i14)}

\begin{itemize}[nosep]
\item SQMS Center received \$125M for 5-year renewal (November 2025).
\item SQMS has achieved $Q_L \sim 10^{10}$ in 1.3~GHz SRF cavities
  for the Dark SRF dark photon search.
\item The Dark SRF exclusion ($\chi < 1.6\ee{-16}$ for dark photons)
  does \emph{not} constrain DFD's scalar coupling (W4i12: different physics).
\item A March 2026 Fermilab study advances quantum sensors for
  particle physics and dark matter detection, confirming continued
  institutional investment.
\end{itemize}

\section{Protocol Summary (Unchanged from W4i10/W4i13)}

\begin{enumerate}[label=\textbf{Phase I.\arabic*:}]
\item Procure/fabricate 4 TESLA-type 1.3~GHz Nb cavities (RRR $\geq$ 300,
  N-doped). Two at 1.3~GHz, two at 2.6~GHz (second harmonic).

\item Achieve baseline $Q_0 = 5\ee{10}$ at 1.3~GHz (demonstrated by SQMS).
  Target: $Q_0 = 2\ee{11}$ with advanced surface treatments.

\item Measure Q-ratio: $R_Q \equiv Q_0(2.6\,\text{GHz}) / Q_0(1.3\,\text{GHz})$.
  \begin{itemize}[nosep]
  \item DFD prediction: $R_Q = 0.49$ (psi-channel loss scales as $\omega^2$).
  \item BCS prediction: $R_Q = 3.41$ (BCS surface resistance scales as
    $\omega^{-2} e^{-\Delta/k_BT}$).
  \end{itemize}

\item Implement Fock-state enhancement: $\langle n | a^\dagger a | n \rangle = n$
  provides $n$-fold sensitivity boost for psi-channel loss measurement.
  At $n = 10$ (demonstrated in microwave cavities): effective sensitivity
  $\lbone < 7.8\ee{-10}/\sqrt{10} = 2.5\ee{-10}$.

\item Triple-blind protocol: TLS mitigated via N-doping; nonequilibrium
  quasiparticle (NEQ) mitigated via magnetic shielding + phonon traps.
\end{enumerate}

\section{Inconsistency Flag}

\inconsistency{The derivation of the psi-channel loss rate
$\Gamma_\psi \propto |\lambda-1|^2 / \Qpsi$ from the DFD action
principle (section02 Eq.~2.6) has not been independently verified
outside the DFD research programme. The SQMS bound $\lbone < 1.56\ee{-9}$
(W3i3) depends entirely on this derivation. If the loss rate formula
is incorrect by a factor of $X$, the bound weakens to
$\lbone < 1.56\ee{-9} \times \sqrt{X}$. This should be flagged to
Agent 1 (Formalism) for independent re-derivation.}


%=============================================================================
%=============================================================================
\part{Protocol 5: Si$_3$N$_4$ MEMS In-Gap Test}
\label{part:mems}
%=============================================================================
%=============================================================================

\begin{tcolorbox}[colback=blue!3!white, colframe=proposalcolor,
  title=\textbf{PROTOCOL 5: MEMS IN-GAP $\Qpsi$ THRESHOLD TEST}]
\textbf{Budget}: \$200--500K\\
\textbf{Timeline}: 12--18 months\\
\textbf{Suppliers}: Norcada (Alberta), Atomica (Santa Barbara),
  university cleanrooms\\
\textbf{Personnel}: 1 PI (optomechanics group) + 1 postdoc\\
\textbf{DFD prediction}: $\Qpsi$ threshold $2.2\ee{8}$ (4.5$\times$
  margin at $Q_{\rm mech} = 10^9$)
\end{tcolorbox}

\section{Key Specifications}

\begin{center}
\renewcommand{\arraystretch}{1.2}
\begin{tabular}{@{}lll@{}}
\toprule
\textbf{Parameter} & \textbf{Required} & \textbf{Demonstrated} \\
\midrule
Material & High-stress Si$_3$N$_4$ & Commercial (Norcada) \\
$Q_{\rm mech}$ & $\geq 10^9$ & $6.8\ee{6}$ (1D, mK) \\
Temperature & $\leq 20$~mK & Demonstrated \\
Readout & Fiber interferometric & Standard \\
Frequency & 100~kHz--1~MHz & Tunable by geometry \\
\bottomrule
\end{tabular}
\end{center}

\caution{The demonstrated $Q_{\rm mech}$ of $6.8\ee{6}$ at cryogenic
temperatures is still $\sim 150\times$ below the $10^9$ requirement.
The path to $10^9$ relies on soft-clamped ``trampoline'' geometries
(demonstrated at room temperature with $Q \cdot f > 10^{13}$~Hz) being
maintained at millikelvin temperatures. This has not been demonstrated
and represents the primary technical risk.}

\section{DFD Observable}

The ``in-gap'' test exploits the prediction that the psi-field provides
an additional dissipation channel with quality factor $\Qpsi$. If
$Q_{\rm mech} > \Qpsi$, the mechanical $Q$ will be limited by
$\Qpsi$ rather than by intrinsic material losses. The DFD prediction
is $\Qpsi = 2.2\ee{8}$ at the SQMS bound.

\textbf{Test}: Cool a Si$_3$N$_4$ MEMS resonator to 20~mK.
Measure $Q$ vs.\ temperature. If $Q$ saturates at $\sim 2\ee{8}$
despite improved material quality, this is consistent with
psi-channel dissipation. If $Q$ continues to rise above $10^9$,
DFD's $\Qpsi$ prediction is falsified at that level.

\section{Supply Chain (Confirmed W4i10)}

\begin{itemize}[nosep]
\item Norcada Inc.\ (Edmonton, Alberta): Commercial Si$_3$N$_4$ membranes,
  high-stress, custom geometries. Lead time: 4--8 weeks.
\item Atomica Corp.\ (Santa Barbara, CA): Custom MEMS fabrication,
  including soft-clamped designs. Lead time: 8--12 weeks.
\item University cleanrooms: Stanford, Yale, TU Delft for
  specialized geometries not available commercially.
\end{itemize}


%=============================================================================
%=============================================================================
\part{Protocol 6: DFD Boltzmann Code (\texttt{mochi\_class})}
\label{part:boltzmann}
%=============================================================================
%=============================================================================

\begin{tcolorbox}[colback=red!5!white, colframe=warncolor,
  title=\textbf{URGENT: DFD BOLTZMANN CODE --- REVISED USING \texttt{mochi\_class}}]
\textbf{Budget}: \$40--80K (revised down from \$80--150K)\\
\textbf{Timeline}: 2--3 months (revised down from 4.5 months)\\
\textbf{Platform}: \texttt{mochi\_class} (arXiv:2407.11968)\\
\textbf{Personnel}: 1 postdoc with Boltzmann code experience\\
\textbf{Deliverable}: DFD predictions for $D_H/r_d$, $S_8(l)$,
  CMB power spectra, neutrino mass constraints
\end{tcolorbox}

\section{Why \texttt{mochi\_class} Supersedes \texttt{hi\_class}}

\begin{enumerate}[nosep]
\item \texttt{mochi\_class} accepts general input functions of time,
  removing the need to hard-code the DFD $\alpha$-functions.
\item Stable basis functions ensure ghost-free and gradient-stable
  evolution --- critical for DFD's nonlinear $\mu$-crossover.
\item Validated against \texttt{hi\_class} and \texttt{EFTCAMB}.
\item Reduces custom code from $\sim$2,520 lines (W4i12 estimate for
  raw \texttt{hi\_class} modification) to $\sim$500--800 lines of
  configuration, wrapper, and post-processing code.
\end{enumerate}

\section{DFD Mapping onto Horndeski $\alpha$-Functions}

From W4i12 Boltzmann Spec:
\begin{align}
\alpha_M(a) &= \frac{1}{aH}\frac{d\psi_0}{d\ln a}, \\
\alpha_B &= 0, \\
\alpha_T &= 0.
\end{align}

The background $\psi_0(a)$ is determined by the cosmological field
equation:
\begin{equation}
\frac{d\psi_0}{d\ln a} = -\frac{4\pi G}{c^2 H^2} \cdot
\frac{(\rho - \bar{\rho}) c^2}{\mu(|\nabla\psi_0|/a_*)},
\end{equation}
where $\mu(x) = x/(1+x)$ and $a_* = 2a_0/c^2$.

\section{Implementation Steps}

\begin{enumerate}[label=\textbf{B.\arabic*:}]
\item Install \texttt{mochi\_class} from GitHub. Verify reproduction of
  $\LCDM$ and standard Horndeski test cases.
\item Implement $\alpha_M(a)$ as a tabulated function using the DFD
  background solution $\psi_0(a)$, computed by integrating the
  cosmological field equation from $a = 10^{-3}$ to $a = 1$.
\item Run: CMB TT/TE/EE power spectra, matter power spectrum $P(k)$,
  lensing convergence $C_l^{\kappa\kappa}$, BAO distances
  ($D_H/r_d$, $D_M/r_d$).
\item Compare with Planck 2018 + DESI DR2 + DES Y5 data.
\item Generate DFD-specific $D_H/r_d$ predictions at DESI redshift bins
  (feeding back into Protocol~1).
\end{enumerate}


%=============================================================================
%=============================================================================
\part{Inconsistencies and Open Issues}
\label{part:inconsistencies}
%=============================================================================
%=============================================================================

\section{Flagged Inconsistencies}

\begin{enumerate}

\item \inconsistency{SQMS bound derivation chain.}
The $\lbone < 1.56\ee{-9}$ bound (W3i3) derives from: DFD action
(Eq.~2.6) $\to$ EM-psi coupling Lagrangian $\to$ psi-phonon vertex
$\to$ cavity loss rate $\Gamma_\psi = |\lambda-1|^2 \omega^2 U_{\rm EM}
/ (c^4 \Qpsi)$ $\to$ comparison with observed $Q_0$. The intermediate
step (psi-phonon vertex) requires showing that the scalar field's
back-reaction on EM modes produces a frequency-dependent loss that
scales as $\omega^2$. This $\omega^2$ scaling is what produces the
Q-ratio diagnostic (0.49 at 2.6/1.3~GHz). If the scaling is $\omega^n$
with $n \neq 2$, the Q-ratio prediction changes and the bound shifts.
\textbf{Action}: Agent 1 should independently derive the psi-channel
loss rate from first principles and verify the $\omega^2$ scaling.

\item \inconsistency{DESI $D_H/r_d$ sign-change location.}
The predicted $z_{\rm cross} = 1.78 \pm 0.12$ (W4i11) was derived from
the analytic psi-screen template. The DESI DR2 Ly$\alpha$ bin is at
$z_{\rm eff} = 2.33$, which is above $z_{\rm cross}$. The QSO bin at
$z_{\rm eff} = 1.49$ is below $z_{\rm cross}$. There is \emph{no DESI
bin directly at $z_{\rm cross}$}. The sign-change test therefore relies
on interpolation between the QSO and Ly$\alpha$ bins. With only 2 data
points bracketing $z_{\rm cross}$, the test may be inconclusive.
\textbf{Mitigation}: DESI DR3 will add additional tracers and finer
binning. Also, the Reboucas et al.\ analysis (arXiv:2510.04179) uses
continuous $D_H/r_d$ ratios rather than binned values, which may
provide better resolution near $z_{\rm cross}$.

\item \inconsistency{Passive Superiority vs.\ technology patent feasibility.}
W4i13 correctly notes this is not a logical contradiction (technology
uses resonant buildup at $C \geq 1$; detection uses passive observation
at $C \ll 1$). However, the technology patents require
$C = \lambda^2 Q_{\rm EM}^2 F_\psi / (M_{\rm eff} \omega^2 / c^4) \geq 1$.
With $M_{\rm eff} = 3.01\ee{6}$~kg (corrected from $10^{-6}$~kg,
W3i4), achieving $C \geq 1$ requires $Q_{\rm EM} \geq 2\ee{11}$ at
$\lbone = 1.56\ee{-9}$. SQMS has demonstrated $Q_0 = 5\ee{10}$ ---
this is $4\times$ below the technology threshold. \textbf{The gap
between current $Q_0$ and technology $C = 1$ is real and should be
honestly communicated.}

\item \inconsistency{MEMS $Q_{\rm mech} = 10^9$ requirement.}
The best demonstrated cryogenic Si$_3$N$_4$ mechanical $Q$ is
$\sim 7\ee{6}$ (1D nanobeam at mK temperatures). The required
$Q_{\rm mech} = 10^9$ is $\sim 150\times$ higher. While room-temperature
$Q \cdot f > 10^{13}$~Hz has been demonstrated (implying $Q > 10^9$
at $f < 10^4$~Hz), maintaining this at millikelvin temperatures
with soft-clamped geometries has \emph{not} been demonstrated.
This is a genuine technical risk that could delay the MEMS experiment.
\end{enumerate}

\section{Items NOT Inconsistent (Confirmation)}

\begin{enumerate}
\item Dark SRF invalidation (W4i12): Confirmed. $G/c^4 = 9.3\ee{-45}$
  kills active EM-to-$\psi$ conversion. SQMS bound is passive (immune).
  No contradiction.

\item EFE screening of wide binaries (W4i11): Confirmed. $a_{\rm ext}
  = 1.8\,a_0$ screens the MOND anomaly to $+10$--12\% at $>10$~kAU.
  Consistent with Banik et al.\ 2024 null result. No contradiction
  with DFD theory.

\item DESI tension at 3.1--3.5$\sigma$ (W4i11): This tension is in
  the \emph{CPL parameterization}, which W4i11 correctly identifies as
  the wrong observable. The $D_H/r_d$ sign-change test (Protocol~1)
  uses the correct DFD-specific observable. No internal contradiction,
  but the tension in CPL remains an honest negative until resolved.
\end{enumerate}


%=============================================================================
%=============================================================================
\part{Consolidated Programme: Updated Budget and Timeline}
\label{part:consolidated}
%=============================================================================
%=============================================================================

\section{Complete Experimental Programme (W4i14)}

\begin{center}
\renewcommand{\arraystretch}{1.3}
\begin{longtable}{@{}clcccl@{}}
\toprule
\textbf{Rank} & \textbf{Experiment} & \textbf{Cost} & \textbf{Time} &
  \textbf{Status} & \textbf{W4i14 Update} \\
\midrule
\endhead
0a & $D_H/r_d$ sign change & \$0 & 2--4 wks & \textbf{EXECUTE NOW} &
  Full protocol, data identified \\
0b & GPS 59~ps archival & \$50--100K & 6 mo & ACTIONABLE &
  Satellites, receivers, software specified \\
0c & LIGO O4 scalar & \$100--500K & 12 mo & ACTIONABLE &
  O4a data available, pipeline specified \\
0d & Rotation curves & \$0 & 3--6 mo & Unchanged & --- \\
\midrule
$\infty$ & DFD Boltzmann code & \$40--80K & 2--3 mo & \textbf{REVISED} &
  \texttt{mochi\_class} reduces effort $2\times$ \\
\midrule
I & SQMS Phase~I & \$3.2M & 24 mo & ACTIONABLE &
  Inconsistency flagged (loss rate derivation) \\
II & Si$_3$N$_4$ MEMS & \$200--500K & 18 mo & Unchanged &
  $Q_{\rm mech}$ gap noted as technical risk \\
\midrule
N1 & Torsion pendulum readout & \$300--600K & 24 mo & Technology dev &
  Repurposed from killed direct force \\
N2 & Atom interferometry & \$50--100K & 6 mo & Feasibility study &
  Edge of 2028--2030 sensitivity \\
\midrule
\multicolumn{6}{l}{\textbf{KILLED (confirmed from W4i13):}} \\
--- & Torsion direct force & --- & --- & KILLED & $G/c^4$: 8 orders short \\
--- & Optical clock altitude & --- & --- & KILLED & Signal $10^{-41}$ \\
--- & Neutron interferometry & --- & --- & KILLED & Signal $10^{-22}$ rad \\
--- & Casimir scalar probe & --- & --- & KILLED & GR-standard, not DFD \\
--- & Muon $g{-}2$ & --- & --- & NULL & $10^{-15}$, publish as null \\
--- & Dark SRF reinterpretation & --- & --- & INVALIDATED & W4i12 \\
\bottomrule
\end{longtable}
\end{center}

\textbf{Total Phase 0}: \$90--680K (archival analyses + Boltzmann code).

\textbf{Total Phase I}: \$3.2M (SQMS, unchanged).

\textbf{Total Phases 0+I}: \$3.3--3.9M over 30 months.


\section{Decision Tree (Updated)}

\begin{center}
\fbox{\parbox{0.92\textwidth}{
\textbf{Step 0 (NOW, Week 1--2)}: Check DESI DR2 paper tables and
Reboucas et al.\ figures for $D_H/r_d$ residual structure. This
takes \emph{hours}, not weeks.\\[4pt]
\textbf{Step 1 (Weeks 2--4)}: Full $D_H/r_d$ sign-change analysis
(Protocol~1). If non-monotonic at $>2\sigma$: fast-track Boltzmann code.\\[4pt]
\textbf{Step 2 (Months 1--3)}: DFD Boltzmann code using
\texttt{mochi\_class} (Protocol~6). Generate rigorous $D_H/r_d$,
$S_8(l)$, CMB predictions.\\[4pt]
\textbf{Step 3 (Months 1--6)}: GPS 59~ps archival (Protocol~2).
Independent of DESI outcome.\\[4pt]
\textbf{Step 4 (Months 3--15)}: LIGO O4 scalar mode (Protocol~3).
Start with O4a; expand as O4b/O4c released.\\[4pt]
\textbf{Step 5 (Months 6--30)}: SQMS Phase~I (Protocol~4).
Proceed regardless of DESI outcome (technology patents independent).\\[4pt]
\textbf{Step 6 (Months 12--30)}: MEMS prototype (Protocol~5).
Parallel with SQMS.\\[4pt]
\textbf{Step 7 (2028+)}: Euclid DR1, CLONETS-DS, JUNO, CMB-S4.
}}
\end{center}


%=============================================================================
\section*{Summary of Derivation Chains}
%=============================================================================

\begin{enumerate}
\item \textbf{$D_H/r_d$ sign change}: Action (section02 Eq.~2.6)
  $\to$ $n = e^\psi$ optical metric (Eq.~2.1) $\to$ $\mu(x) = x/(1+x)$
  crossover $\to$ cosmological $x(z) = a_H(z)/a_0$ $\to$ psi-screen
  modifies $H(z)$ non-monotonically $\to$ $\Delta(D_H/r_d)$ sign change
  at $z_{\rm cross} = 1.78 \pm 0.12$ $\to$ high-$z$ suppression
  matches Reboucas et al.\ 2.3--2.5$\sigma$ trend.
  \textbf{Testable NOW with published data.}

\item \textbf{GPS 59~ps}: Action (Eq.~2.6) $\to$ Sun-sourced
  $\psi_{\rm grav}$ gradient $\to$ nonreciprocal timing
  $\delta t = \Delta\psi / (2c)$ $\to$ 59~ps at GPS altitude $\to$
  multi-GNSS ratio $R = 1.070$ (altitude-dependent, atmospheric-independent).
  \textbf{Testable with IGS archival data and PRIDE PPP-AR.}

\item \textbf{LIGO 6th channel}: Action (Eq.~2.6) $\to$ $\Box\psi$
  sourced by merger $\to$ scalar breathing mode $h_b$ $\to$
  common-mode rejected in differential LIGO $\to$ reach $1.5\ee{-14}$
  (O4 squeezing). \textbf{O4a strain data on GWOSC now.}

\item \textbf{SQMS Q-ratio}: Action (Eq.~2.6) $\to$ EM-psi coupling
  $\to$ psi-channel loss $\Gamma_\psi \propto \omega^2$ $\to$
  $R_Q = 0.49$ (DFD) vs.\ $3.41$ (BCS). \textbf{Requires independent
  verification of $\omega^2$ scaling (flagged).}

\item \textbf{MEMS in-gap}: Action (Eq.~2.6) $\to$ psi-phonon coupling
  $\to$ $\Qpsi = 2.2\ee{8}$ $\to$ mechanical $Q$ saturates if
  $Q_{\rm mech} > \Qpsi$ $\to$ observable plateau.
  \textbf{Requires $Q_{\rm mech} = 10^9$ not yet demonstrated
  cryogenically (technical risk).}
\end{enumerate}


\end{document}
